\newacronym{ia}{IA}{Inteligencia Artificial}
\newacronym{csv}{CSV}{Comma Separated Values}
\newacronym{ai}{AI}{Artificial Intelligence}
\newacronym{api}{API}{Application Programming Interface}
\newacronym{RNN}{RNN}{Recurrent Neural Network}
\newacronym{LSTM}{LSTM}{Long Short-Term Memory}
\newacronym{npc}{NPC}{Non Playable Character}
\newacronym{rpg}{RPG}{Rol-Playing Game}
\newglossaryentry{BT}{
    name=Behavior Tree,
    description={Un Behavior Tree es un modelo de ejecución de tareas de forma jerárquica basado en árboles usada en el campo de la robótica y la inteligencia artificial}
}

\newglossaryentry{BB}{
    name=Blackboard,
    description={La Blackboard es la memoria compartida usada en un Behavior Tree. Ésta funciona como un diccionario de pares clave-valor, que pueden ser variables de entrada/salida.}
}

\newglossaryentry{motor}{
    name=motor de videojuegos,
    description={Un motor de videojuegos es un software compuesto por varias librerías (motor de renderizado, motor de físicas, motor de audio, etc.) que permiten el desarrollo y la ejecución de un videojuego.}
}
\newglossaryentry{LLM}{
    name=LLM,
    description={Un LLM (Large Lenguage Model) es un modelo de aprendizaje automático profundo (deep learning) que ha sido entrenado con grandes cantidades de texto con propósito general de prediccónn de palabras.}
}

\newglossaryentry{Singleton}{
    name = Singleton,
    description = {El patrón Singleton es un patrón de programación que consiste en garantizar que una clase tenga una instancia única y que se pueda acceder globalmente a ésta.}
}

\newglossaryentry{Abstract Factory}{
    name = Abstract Factory,
    description = {El patrón Abstract Factory es un patrón de programación que consiste en mantener una clase que permita crear instancias de objetos relacionados gracias al polimorfismo.}
}

\newglossaryentry{Ollama}{
    name = Ollama,
    description = {Ollama es una aplicación basada en el motor de inferencia ``llama.cpp'' que permite descargar y ejecutar modelos localmente.}
}

\newglossaryentry{Fine-tuning}{
    name = Fine-tuning,
    description = {El Fine-tuning es una técnica de Machine Learning que consiste en reentrenar ciertas partes de modelos ya preentrenados con nuevos datos con el objetivo de especializarlo en un tema concreto.}
}

\newglossaryentry{RAG}{
    name = RAG,
    description = {RAG (Retrieval Augmented Generation) es una técnica de Inteligencia Artificial que permite mejorar las respuestas de modelos de lenguaje buscando información en documentos externos que tengan que ver con la pregunta del usuario.}    
}

\newglossaryentry{Embedding}{
    name = Embedding,
    description={Un embedding es un medio para representar la información compleja de forma numérica. Esto se consigue transladando la información a un espacio vectorial en el que a cada unidad se le determina unas coordenadas según el contexto que tiene a su alrededor. Así, conceptos cercanos se mantienen matemáticamente cerca y son más sencillos de relacionar.}
}

\newglossaryentry{Almacén vectorial}{
    name = Almacén vectorial,
    description = {Un almacén vectorial (o vector store) es una base de datos especializada en el almacenaje de la información en forma de vectores (o embeddings), la cual permite recuperar los datos por similitud matemática a otros vectores.}
}

\newglossaryentry{FAISS}{
    name = FAISS,
    description = {FAISS (Facebook AI Similarity Search) es una libería de Meta que permite la búsqueda y el agrupamiento eficiente de vectores densos con el fin de facilitar la búsqueda por similitud.}
}
